
\AddSymbolsB 
\pagecolor{cyan!10!gray}
\color{white}
\quad\quad \lettrine[,lraise=0.2,lhang=0.5]{\textbf{伊}}{藤}佐一百无聊赖地躺在床上，玩弄着奖章。他的床头开着一盏台灯，照向了别处。

他突然害怕了起来。

他害怕自己可能到死都无法理解这枚奖章出现的含义。

可是他一想到昨天教授明确表示的对自己的欣赏，便安下了心。

\twkkai{教授一定会积极地帮我答疑解惑，或者……弄清真相的；}伊藤佐一关掉了台灯，睡着了。

\begin{center}
——————
\end{center}

第二天，伊藤佐一按时赶到了课室。课室里早就坐满了人，但中间仍有一两排稀疏的座位。伊藤佐一坐在了其中一个两边都没有人的座位上。

他回头张望着。铃木和子似乎没有来。

他转而看向手表——上课已经3分钟了；前门早已关闭。

他再次回头张望着，扫视着一排又一排的人群。平时连别人眼睛都不敢对视的他，这次冒着被集体注视的风险，一个一个地找。

还是没有。没有她的身影。

此时的他，感觉到的不是疑惑，是焦灼。他又把手揣进口袋里，拿出那枚奖章，仔细看了起来。

奖章依然是崭新的，但刻的东西的的确确已经和昨天的不同了：
\[
\forall a,d\in\mathbb{N},\ \gcd(a,d)=1\ \Rightarrow\ \#\{p\in\mathbb{P}:p\equiv a\pmod d\}=\infty
\]

教授依旧在台下，讲回了教材内容，但就是和奖章上的铭文风马牛不相及。

他看了一眼课室的钟。离下课还有十七分钟。他盯着它，出了神，还隐隐约约听到有人呼喊他的名字。

而今天，既没有猫的出现，也没有教授的激情讲解；再加上铃木和子的缺席，伊藤佐一感到了学习生涯中史无前例的无聊。

……

下课铃声响起，伊藤佐一飞也似地第一个跑到教授跟前。可是没等伊藤佐一出声，教授先发话了。

\twkkai{伊藤同学，你今天的表现不是很好喔——我叫了几次你的名字，你都没有应声……}

伊藤佐一仿佛没有听到，盯着教授，挠了挠头。\twkkai{我……}

他摊开了双手。奖章显露了出来。

\twkkai{又是……那枚奖章？}教授问道。

伊藤佐一点了点头。

教授端详了许久，并不作声。伊藤佐一不知道教授是否和自己一样，为奖章的意义所困；抑或只是专注于铭文的含义。

\twkkai{很好，很好；你的奖章很有趣，可惜我的办公室不是探讨它的最佳场所。伊藤同学，你愿意一同和我出去走走吗？}教授扶了扶眼镜，把奖章还给了伊藤佐一。

伊藤佐一再次点了点头。

两人走出了课室。

……

来到河边，伊藤佐一率先问话。

\twkkai{教授，你有没有听说过班里有个叫铃木和子的同学？}

教授沉思了一会儿。

\twkkai{听说过，不过印象不深；怎么了？}

\twkkai{没事，没事；我和她是要好的同学，可是她今天没来上学；}伊藤佐一撒着谎。说着说着，自己的脸“唰”一下地红了。

\twkkai{没来上学，很正常嘛——可能是身体不舒服；}教授云淡风轻地说。

可是伊藤佐一突然感到全身酥麻。不好的联想在他脑海里涌现。

河边的风吹了过来，拂动着路边的杉树叶。伊藤佐一耳边突然响起一阵轻快的音乐——原来是从岸上的咖啡馆传来的。

伊藤佐一眺望着咖啡馆，想起铃木和子，又想起奖章；便把奖章拿了出来，看着。上面的铭文没有变化，却逐渐变得模糊；这使得伊藤佐一的手心又开始出汗了。铭文变得更为模糊了；伊藤佐一猛地将奖章放回裤袋，腾出双手，把手掌贴在裤子上反复摩擦着。

\twkkai{累了吧，伊藤同学？要不要去喝杯咖啡？}教授扭过头，慈祥地问道。

伊藤佐一不作声，依旧像往常一样点了点头。

他们穿过热闹的自行车流，踏进了咖啡馆门口。香气把伊藤佐一的注意力带回到此时此刻。

\twkkai{你喜欢喝什么？我请你；}教授亲切地微笑着。

\twkkai{我……拿铁吧，谢谢；}伊藤佐一不敢支吾很久——毕竟是第一次和教授面对面正式对话——假装爽快地回答着。

教授去了点餐台；伊藤佐一找了最近的座位坐下，把双肘枕在桌面上；可是桌面冷冰冰的，他又不得不把手放下来。

教授就那么站在台前，什么都不干，静静地等待了十分钟；这让他和忙碌的咖啡师形成了鲜明的对比。伊藤佐一很久没有见过一个人能如此悠闲了；他也的确不喜欢人们忙碌的状态，仿佛他们都陷入了一种理所当然。

外面突然下起了磅礴大雨，噼里啪啦地打在咖啡馆的落地窗上；坐在馆内，想看清河对岸的建筑已经不可能了。可是咖啡馆依旧人声沸腾；伊藤佐一的内心由此而竟然感到无比地平静。他看向咖啡馆内墙上的标语：
\begin{quote}
\begin{center}
\twkkai{暗处不是光的缺席，是光的另一种形式。}
\end{center}
\end{quote}

\twkkai{这句标语会不会和铭文有些许的联系呢？}

伊藤佐一沉思着。此时，教授已经拿着两杯拿铁，朝他走来。

\twkkai{来吧，这是你的；}教授把其中一杯拿铁推向伊藤佐一。

雨停了。已是傍晚，街边的路灯一盏隔着一盏地亮起。

教授用吸管随意地搅拌着拿铁。

\twkkai{伊藤同学，你看看窗外的街灯；}教授用空闲的手指了指窗外。

伊藤佐一看向了那一列交替亮着的街灯，却不明白教授话里的含义。

\twkkai{它们就是你的奖章铭文内容的特例。你玩过跳房子游戏吧；}教授耐心地问道。伊藤佐一点了点头。

\twkkai{假如你有无限的精力，可以从第一盏亮着的灯跳到第三盏亮着的灯，再从第三盏跳到第五盏，以此类推……你会发现脚下用粉笔画的数字，有无穷多是素数……}

\twkkai{而从另一个方面来说，没有暗着的灯，亮着的灯便无从为人们所得知了；所以……素数虽是瑰宝，但仍要经由合数的衬托才得以放出光芒啊……}

教授用手指了指身后那句标语。

伊藤佐一极为专注地听着，为教授话语中展现的魔力所吸引，而并没有留意到馆外的一只猫——它正在离咖啡馆门口不远处的一棵树下徘徊。此时，它一跃而起，跳到了一盏没有亮起的街灯顶端。那盏灯突然闪了闪，接着便稳定地发出耀眼的光芒；而所有处于熄灭状态的街灯，此时此刻同时亮起。

而此时，教授从口袋里掏出了一张折了多次的草稿纸——展开后竟有整张桌面那么大。

\twkkai{……而你的奖章，本质上是要证明……}

教授则停顿了一下，又指了指草稿纸。

\twkkai{……这种“亮”，不是偶然的；我们把它写成数学……}
\begin{gather*}
\text{\twkkai{定义集合\(A\subset P\)的密度为：}}
\delta(A)=\lim_{s\to1^{+}}\frac{\sum_{p\in A}p^{-s}}{\sum_{p\in P} p^{-s}}\\
\text{\twkkai{密度存在且为正\(\Rightarrow\)无穷多素数.}}
\end{gather*}\index{mi@密度 $\delta(A)$}\index{mi@密度 $\delta(A)$!定义式：$\lim_{s\to1^+}\frac{\sum_{p\in A}p^{-s}}{\sum_{p\in P}p^{-s}}$}\index{mi@密度 $\delta(A)$!密度为正蕴含无穷多素数}
\ \ \quad \twkkai{那么，要得到这个结论，我们这次的终极武器，就和1837年狄利克雷发明的\(L\)-函数有关了；你也可以说它是被发现的——不过这就是数学哲学家的研究课题了……}

\twkkai{\(L\)-函数？1837？那不是昨天……}

\twkkai{是的——就是我们昨天讲过的内容；}教授点了点头——他显然没有听懂伊藤佐一的言外之意。

墨水继续从教授拿着的笔中挤出，附着在纸面上。\twkkai{看到这个密度的定义式，我们自然想要求出分子分母的渐近式是什么，对吧；}教授喃喃自语着。

\twkkai{自然，我们这里还要介绍\(L\)-函数的近亲——黎曼\index{li@黎曼（B. Riemann）}\(\zeta\)函数\index{li@黎曼（B. Riemann）!黎曼 $\zeta$ 函数 $\zeta(s)$}。还有，让我告诉你吧——其实狄利克雷和黎曼也是师生关系\index{dilikelei@狄利克雷（P.\,G.\,L.\ Dirichlet）!与黎曼的师生关系}，和我们一样；}教授看着伊藤佐一，打趣地说着。

伊藤佐一显然没有听懂教授的言外之意。他看着教授的笔继续不紧不慢地吐着墨水。
\[\log \zeta(s)=\sum_{p\in P}\log\frac{1}{1-p^{-s}}=\sum_{p\in P}\sum_{n\geq1}\frac{1}{np^{ns}}=\sum_{p\in P}\frac{1}{p^s}+\sum_{p\in P}\sum_{n\geq2}\frac{1}{np^{ns}}\]

 \twkkai{根据这个运算，数学家们很久以前便得到了一个有趣的结果——}
\begin{gather*}
\text{\twkkai{对属于收敛域内的实数\(s\), 当\(s\to1^+\)时, }}\sum_{p\in P}\frac{1}{p^s}\sim\log\frac{1}{s-1}.
\end{gather*}

门外却突然无缘无故传来一阵欢呼声——原来那只猫正腾空飞跃在街灯之间；街灯下早已围起了一群人。

而教授依旧专注地讲解着；但伊藤佐一已经分心了——他逐渐开始担忧起来，希望教授讲快点。可是教授仍是那么慢条斯理。

伊藤佐一本来有许多问题想问教授，以此将放在猫身上的注意力夺回来；可是他此时已经紧张得忘光了。他时常瞟向窗外。

这时，咖啡馆的门被其中一位顾客推开；猫趁着这个间隙，偷偷地钻了进来。

伊藤佐一转而死死盯着那只猫；而那只猫也紧紧地盯着他，一动也不动。但时而有顾客路过，伸出手指勾了猫的下巴几次，猫便十分不耐烦，直接跳到了伊藤佐一座位旁。

伊藤佐一故意扭头再次看向教授，假装不在乎猫。他像是在专注地思考，实则只是盯着教授的手指发呆。猫见状后，便安定地盘坐了下来，开始舔自己的后脚掌。

教授全程都没有留意到伊藤佐一和猫的舞台剧。

\twkkai{那么……分母已经求出来了，分子\(\sum_{p\in A}p^{-s}\)的渐近式应该怎么得到呢？}

教授终于又出声了；他的头发优雅地侧向一边，轻轻地搭在额头上。他的写字的手已经分布着些许皱纹，可是并不颤抖——反而极为平稳地书写着。

写着写着，教授突然看向伊藤佐一。

\twkkai{你是想我卖关子好，还是不卖关子好呢？}

\twkkai{喵……}

猫突然把其中一只前脚掌搭在了伊藤佐一大腿上；它似乎闻到了什么。

伊藤佐一十分纳闷——教授不应该没有留意到猫。\twkkai{可能是过于专注了吧；}他想着。

他犹疑地盯着猫，又看回教授，点了点头。

\twkkai{好！既然如此……我们需要用到两条引理……}

猫看向教授；而它的猫爪却把伊藤佐一的大腿抓得越来越紧了。

\twkkai{第一条说的是，对于模 \(m\) 的特征标 \(\chi\)，当实数 \(s\) 从右侧趋近于 \(1\) 时：如果 \(\chi\) 是主特征标（即 \(\chi = 1\)），那么对所有不被 \(m\) 整除的素数 \(p\) 求和 \(\frac{\chi(p)}{p^s}\)，这个和趋于无穷大——$$\sum_{\substack{p \in \mathbb{P} \\ p \nmid m}} \frac{\chi(p)}{p^s} \sim \log \frac{1}{s-1}, \quad s \to 1^+ ；$$如果 \(\chi\) 是非主特征标，那么这个和是有界的，也就是说它不会发散，而是趋近于某个有限的值。我们把它写下来……$$\sum_{\substack{p \in \mathbb{P} \\ p \nmid m}} \frac{\chi(p)}{p^s} = O(1), \quad s \to 1^+$$}\index{te@特征标$\chi$!主特征标（$\chi=1$）}\index{te@特征标$\chi$!非主特征标（$\chi\neq 1$）}

教授重新拿起了笔，马不停蹄地写着。

此时此刻，伊藤佐一的紧张则到达了顶峰——那只猫正用猫爪死死拽住伊藤佐一手里的奖章。

\twkkai{喵……}

\twkkai{那么，第二条说的是，对于模 \(m\) 的一个剩余类 \(a\)（即素数 \(p\) 要满足 \(p \equiv a \pmod{m}\)），  
所有这种素数的 \(p^{-s}\) 的和，  
可以拆成对所有模 \(m\) 的特征 \(\chi\) 求和。}

\twkkai{简单说，就是……$$\sum_{p \in \mathbb{P}_a} p^{-s} = \frac{1}{\phi(m)} \sum_{\chi \in \hat{G}} \overline{\chi(a)} \left( \sum_{\substack{p \in \mathbb{P} \\ p \nmid m}} \frac{\chi(p)}{p^s} \right).$$}

\twkkai{这条公式说的是——先对每个特征 \(\chi\)，算出所有不被 \(m\) 整除的素数对应的 \(\chi(p)/p^s\) 的和，  
再乘上 \(\overline{\chi(a)}\)——也就是 \(\chi(a)\) 的共轭；  
把这些结果加起来，最后除以 \(\phi(m)\)，  
就正好等于左边那个只属于剩余类 \(a\) 的素数倒数和。}

\twkkai{那么，\(\sum_{p \in P_a} p^{-s} \sim \frac{1}{\phi(m)} \log \frac{1}{s-1}\)这一条结论就呼之欲出了……}

\twkkai{哦，对了，忘了告诉你——\(\phi(m)\)是大名鼎鼎的欧拉函数。}\index{ou@欧拉（L.\ Euler）!欧拉函数 $\phi(m)$}

教授用笔点了点纸面，以示写毕。

\twkkai{喵！}

猫却突然受到惊吓，叼走了伊藤佐一手里的奖章，径直跳到了咖啡桌上，望着教授，呜咽着；随后一跃而起，疯狂跑出了咖啡馆。

而教授终于抬起了头；伊藤佐一早已不见踪影。

教授左右张望着——原来伊藤佐一已经跑远了。教授笑了笑。

\twkkai{这狄利克雷，也挺会挑人……}

接着，他给伊藤佐一发了一条短信。

\twkkai{奖章的事，下次再聊吧。}

教授端起咖啡，啜饮了一口，之后便站了起来，缓缓地走向柜台，帮伊藤佐一付了咖啡钱，接着缓缓地推开门，迎接风的拥抱。

\begin{center}
——————
\end{center}

伊藤佐一追随着猫，疯狂奔跑在雨后撒满粉紫色花瓣的路面上，喘得上气不接下气。烈日之下，伊藤佐一开始出汗，视线开始模糊。

猫似乎终于停了下来。而猫面前似乎坐着一个人。

伊藤佐一也停了下来。此时，他才惊奇地发现，那个人，竟是铃木和子。

\twkkai{你……怎么会在这里？}伊藤佐一首先发话，喘着气，轻声问道。

\twkkai{我出来走走……}

\twkkai{我们昨天约好的……}

\twkkai{非常抱歉……我今天头痛……所以没来……}

\twkkai{这样吗？原来如此；很抱歉我刚才的语气……}

\twkkai{没事的；来，还给你……}

铃木和子似乎知道猫嘴里叼着的东西。趴在膝盖上的猫，自觉地吐出了那块奖章，放在铃木和子的手掌上。她把它递给了伊藤佐一。

伊藤佐一接过奖章。奖章再次变得崭新；而铭文的内容也再一次发生了变化。

\twkkai{我应不应该把这件事告诉她？但是，我知道的是否已经足够……}

伊藤佐一谢过铃木和子，转过身沉思着。

他又转过身。

\twkkai{铃木同学，我有话要对你说……}

猫警觉地竖起耳朵。

\twkkai{……我想知道，你的头痛是什么时候开始的；}他没有意识到，这一句话，让他的脸再次“唰”地红了起来。

\twkkai{就是昨天和你告别后大概二十分钟左右，就开始了；}铃木和子平静地说着。

可伊藤佐一预感到，这平静的背后，躲藏着不祥的征兆。

\twkkai{和你没有关系啦……谢谢你的关心；}铃木和子补充道。

此刻的伊藤佐一，能够清晰地听到自己的心跳声。他不知道如何应答，但他必须说点什么。

\twkkai{嗯……注意要好好休息；我也要谢谢你……那，我先走了；}伊藤佐一和铃木和子招了招手，便转身急步走了起来。猫并没有跟随。

……

回到咖啡馆，伊藤佐一发现教授早已离开；而桌面上留着一张纸。

\newpage
\color{black}
\pagecolor{white}
\twkkai{伊藤同学，鉴于你感兴趣，这是我刚才为你补充的两条定理的证明。你若是有疑问，欢迎明天来办公室问我。切忌在课堂上拿出并提问！（P.S.：之后你若有兴趣，我可再次通过这个定理和你一起探索类域论的美妙花园。）}
\begin{lmm}
 \twkkai{对任何模\(m\)特征标\(\chi\)和实数\(s\to1^+\), 有：}\index{te@特征标$\chi$!模 $m$ 特征}
\begin{gather*}
    \sum_{\substack{p \in P \\ p \nmid m}} \frac{\chi(p)}{p^{s}} \sim \log \frac{1}{s-1} \quad \text{\twkkai{对于}} \chi = 1.
\end{gather*}
\twkkai{左边的级数对于\(\chi\neq1\)是有界的.}
\end{lmm}
\begin{proof}
\twkkai{首先注意到此级数对于 \(s>1\) 是收敛的. 对于 \(\chi=1\) , 级数缺少的是\(\sum_{p\in P}\frac{1}{p^s}\)中的有限多项. 根据结论}
\[\sum_{p\in P}\frac{1}{p^s}\sim\log\frac{1}{s-1},\]
\twkkai{缺少有限项不影响结论的成立, 因此已经证明.}
\twkkai{现假设\(\chi\neq1.\) 对\(L\)-函数取对数可以得到：}
\begin{align*}
\log L(s,\chi)&= \sum_{p \in P} \log \frac{1}{1 - \chi(p) p^{-s}}= \sum_{p \in P} \sum_{n \ge 1} \frac{\chi(p)^n}{n p^{n s}}\\
&= \sum_{p \in P} \frac{\chi(p)}{p^s}
      + \sum_{p \in P} \sum_{n \ge 2} \frac{\chi(p^n)}{n p^{n s}}
\end{align*}
\twkkai{第二个和式\(\sum_{p \in P} \sum_{n \ge 2} \frac{\chi(p^n)}{n p^{n s}}\)在\(s\to1^+\)时是有界的. 由于 \(L(s, \chi)\) 在 \(s = 1\) 的邻域内不为零, 故\(\log L(s, \chi)\) 在此邻域内也是有界的. 因此\(\sum_{\substack{p \in P \\ p \nmid m}} \frac{\chi(p)}{p^{s}}\)在\(s\to1^+\) 时也是有界的.}
\end{proof}
\begin{lmm}
\begin{gather*}
    \sum_{p \in P_{a}} p^{-s}
    = \frac{1}{\phi(m)} 
      \sum_{\chi \in \widehat{G}} \overline{\chi(a)}
      \left( \sum_{\substack{p \in P,\ p \nmid m}} \frac{\chi(p)}{p^{s}} \right)
\end{gather*}
\end{lmm}
\begin{proof}
\twkkai{在证明的最后一步，我们需要将“素数属于 \(P_a\)”这一条件表达为特征和. 关键公式：}
\[
\sum_{\chi\in\widehat{G}} \chi(a)^{-1}\chi(p) = \begin{cases}
\phi(m), & p\equiv a\pmod m,\\
0, & \text{\twkkai{否则}}.
\end{cases}
\]

\twkkai{(考虑群 \(G\)，元素 \(a^{-1}p\bmod m\) 为1时和为 \(\phi(m)\)，否则为0.)}

\twkkai{当 \(s \to 1^+\) 时，只有 \(\chi = 1\) 项贡献渐近 \(\log \frac{1}{s-1}\)，其余项有界。}

\twkkai{对 \(\chi = 1\)：}
  \[
  \sum_{p\in P,\ p \nmid m} \frac{\chi(p)}{p^s} \sim \log \frac{1}{s-1}
  \]

\twkkai{对 \(\chi \neq 1\)：该和有界（因为 \(\log L(s,\chi)\) 在 \(s=1\) 附近有界，且 \(L(1,\chi) \neq 0\)）.}
\end{proof}

\noindent\textbf{结论：}

\twkkai{分子渐近：\(\frac{1}{\phi(m)} \log \frac{1}{s-1}\)}

\twkkai{分母渐近：\(\log \frac{1}{s-1}\)}

\twkkai{因此：}
  \(
  \delta(P_a) = \frac{1}{\phi(m)}\Rightarrow
  \)
  \twkkai{密度为正 \(⇒\) 有无穷多个素数在 \(P_a\) 中.}
\clearpage
\pagecolor{cyan!10!gray}
\color{white}

伊藤佐一倍感失落——他冥冥之中觉察到，似乎假如没有教授和铃木和子，他就要永远独自面对这个不可解的世界了。但他还是拿起了那张沾满咖啡渍的纸。

\twkkai{“切忌在课堂上拿出来并提问……”，为什么呢？}

伊藤佐一抬头观察着。周围的人依旧各做各的事，无暇他顾。

他便把纸整齐地对折，连同奖章塞进裤袋，竭力不去回想今天发生的所有事情——包括数学。


